\begin{lemma}[Second variation of the localizing product]
Let $E$ be a real inner-product space, let $B:E\to\R$ be of class $C^2$ on $E$,
let $f,f_1:\R\to\R$, let $c,u\in E$ and let $r,s,\gamma_2\in\R$. Assume that
\[
  f\ \text{is differentiable at }\rho\ \text{with}\ f'(\rho)=f_1(\rho)
  \quad\text{for every }\rho\ \text{in some neighbourhood of }r,
\]
and that $f_1$ is differentiable at $r$ with $f_1'(r)=\gamma_2$. Then there is a
function $G:\R\to\R$ such that
\begin{enumerate}
\item for every $\tau$ in some neighbourhood of $0$ the function
  \[
    t\longmapsto B(c+tu)\,f(r+ts)
  \]
  is differentiable at $\tau$ with derivative $G(\tau)$, and
\item $G$ is differentiable at $0$ with
  \[
    G'(0)=D^2B(c)[u,u]\,f(r)+2\,DB(c)[u]\,f_1(r)\,s+B(c)\,\gamma_2\,s^2 .
  \]
\end{enumerate}
In other words, along the motion in which the center moves from $c$ with
velocity $u$ and the radius moves affinely from $r$ with velocity $s$, the
product $B(c)f(r)$ has a second derivative at the base point, and that second
derivative is the displayed Leibniz sum.
\end{lemma}
\begin{proof}
Write $\beta(\tau)=B(c+\tau u)$ and $\gamma(\tau)=f(r+\tau s)$, so that the
function in (1) is $\beta\gamma$.

\emph{Step 1: the two straight motions.} The map $t\mapsto c+tu$ is affine, hence
differentiable at every $\tau$ with derivative $u$; likewise $t\mapsto r+ts$ is
differentiable at every $\tau$ with derivative $s$. Both maps send $0$ to $c$
and to $r$ respectively.

\emph{Step 2: $\beta$ is differentiable everywhere, with
$\beta'(\tau)=DB(c+\tau u)[u]$.} Since $B$ is $C^2$ on $E$ it is in particular
differentiable at every point, so $B$ is differentiable at $c+\tau u$ with
derivative $DB(c+\tau u)$. Composing with the motion of Step 1, the chain rule
gives
\begin{equation}\label{eq:ppsv-beta1}
  \beta'(\tau)=DB(c+\tau u)[u]\qquad(\tau\in\R);
\end{equation}
in particular $\beta(0)=B(c)$ and $\beta'(0)=DB(c)[u]$.

\emph{Step 3: $\beta'$ is differentiable at $0$, with
$\beta''(0)=D^2B(c)[u,u]$.} As $B$ is $C^2$, the map $x\mapsto DB(x)$ is $C^1$
on $E$, hence differentiable at $c$, with derivative $D^2B(c)$, a continuous
linear map from $E$ to the space of continuous linear functionals on $E$.
Composing with the motion $t\mapsto c+tu$, which has derivative $u$ at $0$ and
sends $0$ to $c$, gives that $\tau\mapsto DB(c+\tau u)$ is differentiable at $0$
with derivative $D^2B(c)[u]$. Evaluation at $u$ is a continuous linear
functional on the space of continuous linear functionals, so composing once more
and using \eqref{eq:ppsv-beta1},
\begin{equation}\label{eq:ppsv-beta2}
  \beta''(0)=D^2B(c)[u,u].
\end{equation}

\emph{Step 4: $\gamma$ near $0$, and $\gamma'$ at $0$.} The map
$\tau\mapsto r+\tau s$ is continuous and sends $0$ to $r$, so it maps some
neighbourhood of $0$ into the neighbourhood of $r$ on which $f$ is
differentiable with derivative $f_1$. For $\tau$ in that neighbourhood of $0$,
$f$ is differentiable at $r+\tau s$ with derivative $f_1(r+\tau s)$, and the
chain rule with Step 1 gives that $\gamma$ is differentiable at $\tau$ with
\begin{equation}\label{eq:ppsv-gamma1}
  \gamma'(\tau)=f_1(r+\tau s)\,s .
\end{equation}
Taking $\tau=0$: $\gamma(0)=f(r)$ and $\gamma'(0)=f_1(r)s$. Moreover $f_1$ is
differentiable at $r$ with derivative $\gamma_2$, so composing with the motion
once more,
\begin{equation}\label{eq:ppsv-gamma2}
  \tau\mapsto f_1(r+\tau s)s\ \text{is differentiable at }0
  \ \text{with derivative }\gamma_2 s\cdot s=\gamma_2s^2 .
\end{equation}

\emph{Step 5: definition of $G$ and proof of (1).} Put
\[
  G(\tau)=DB(c+\tau u)[u]\,f(r+\tau s)+B(c+\tau u)\,f_1(r+\tau s)\,s
   =\beta'(\tau)\gamma(\tau)+\beta(\tau)\gamma'(\tau)
\]
for $\tau$ in the neighbourhood of $0$ supplied by Step 4, extended arbitrarily
(say by $0$) elsewhere. For such $\tau$, $\beta$ is differentiable at $\tau$ by
Step 2 and $\gamma$ is differentiable at $\tau$ by Step 4, so the product rule
gives that $\beta\gamma$ is differentiable at $\tau$ with derivative
$\beta'(\tau)\gamma(\tau)+\beta(\tau)\gamma'(\tau)=G(\tau)$. This is (1).

\emph{Step 6: proof of (2).} On that neighbourhood, $G=G_1+G_2$ with
$G_1(\tau)=\beta'(\tau)\gamma(\tau)$ and $G_2(\tau)=\beta(\tau)\gamma'(\tau)$.
By Step 3 the factor $\beta'$ is differentiable at $0$ with derivative
$\beta''(0)$, and by Step 4 the factor $\gamma$ is differentiable at $0$ with
derivative $\gamma'(0)$; the product rule gives that $G_1$ is differentiable at
$0$ with
\[
  G_1'(0)=\beta''(0)\gamma(0)+\beta'(0)\gamma'(0).
\]
By Step 2 the factor $\beta$ is differentiable at $0$ with derivative
$\beta'(0)$, and by \eqref{eq:ppsv-gamma2} the factor $\gamma'$ is
differentiable at $0$ with derivative $\gamma_2s^2$; the product rule gives that
$G_2$ is differentiable at $0$ with
\[
  G_2'(0)=\beta'(0)\gamma'(0)+\beta(0)\,\gamma_2s^2 .
\]
Adding, $G$ is differentiable at $0$ with
\[
  G'(0)=\beta''(0)\gamma(0)+2\beta'(0)\gamma'(0)+\beta(0)\,\gamma_2s^2
   =D^2B(c)[u,u]\,f(r)+2\,DB(c)[u]\,f_1(r)\,s+B(c)\,\gamma_2\,s^2,
\]
where the values $\beta(0)=B(c)$, $\beta'(0)=DB(c)[u]$,
$\beta''(0)=D^2B(c)[u,u]$, $\gamma(0)=f(r)$ and $\gamma'(0)=f_1(r)s$ were
substituted from Steps 2, 3 and 4. This is (2).
\end{proof}
